\documentclass[10pt]{article}

\usepackage[T1]{fontenc}
\usepackage{lmodern}
\usepackage[letterpaper,margin=0.85in]{geometry}
\usepackage{booktabs}
\usepackage{enumitem}
\usepackage{microtype}
\usepackage{xcolor}
\usepackage[hidelinks]{hyperref}

\setlist[itemize]{leftmargin=*,itemsep=2pt,topsep=3pt}
\setlist[enumerate]{leftmargin=*,itemsep=4pt,topsep=3pt}
\setlength{\parindent}{0pt}
\setlength{\parskip}{4pt}

\newcommand{\paperTitle}{Event-FedAvg: Stateful Event-Driven Communication for Federated Learning}
\newcommand{\decision}{Weak reject in the current form}

\begin{document}

\begin{center}
  {\Large\bfseries Independent Mock Reviewer Report}\par
  \vspace{4pt}
  {\large IJCNN 2027 submission stress test}\par
\end{center}

\textbf{Manuscript:} \paperTitle

\textbf{Version reviewed:}
repository commit \texttt{ee29e78b9a019b307cd2a4afbc366d5432ded725},
reviewed on 6 September 2026.

\textbf{Status of this document:} This is an internal mock review written from
the perspective of a technically engaged but skeptical IJCNN reviewer. It is
not an official conference review. The score labels below are a transparent
mock scale because the final IJCNN 2027 reviewer form was not available at the
time of review.

\section*{Recommendation and scores}

\begin{center}
\begin{tabular}{@{}p{0.43\textwidth}p{0.18\textwidth}p{0.28\textwidth}@{}}
\toprule
Criterion & Score & Assessment \\
\midrule
Relevance to IJCNN & 4/5 & Clearly in scope \\
Originality & 3/5 & Distinct operator, incremental ingredients \\
Technical soundness & 3/5 & Careful but conditional theory \\
Empirical validation & 2/5 & Promising, presently too narrow \\
Clarity and organization & 4/5 & Strong draft with local ambiguities \\
Reproducibility & 4/5 & Unusually good artifact discipline \\
Overall & 4/10 & \decision \\
Reviewer confidence & 4/5 & High on FL and compression issues \\
\bottomrule
\end{tabular}
\end{center}

The paper is close to the acceptance boundary, but I would not recommend
acceptance on the current evidence. The main reason is not formatting. It is
that the central empirical advantage is estimated from only three held-out
seeds and the experiments do not yet isolate which claimed operator feature
causes that advantage. A focused revision could plausibly move the paper to a
weak accept.

\section*{Summary}

The paper proposes Event-FedAvg, a communication layer around conventional
federated averaging. Each client accumulates weighted local-model deltas in a
leaky coordinate state, emits an address and sign when a threshold is crossed,
fully clears the fired state, and lets the server apply an independently
scheduled model quantum. The protocol includes sparse replay and dense
checkpoint fallback so that downlink synchronization is not omitted from the
traffic calculation.

The paper gives a pathwise event-count bound, an exact decomposition of event
alignment, and a smoothness-based descent inequality under a conditional
alignment assumption. Experiments compare the method with dense FedAvg,
Sign-EF, EF-TopK, and Strom-style threshold pulses on Fashion-MNIST and
CIFAR-10. The reported operating points are communication-efficient within
the evaluated synchronous ten-client setting. A ten-seed factorial audit also
shows how heterogeneity and local-step depth change finite-trajectory
alignment terms.

\section*{Strengths}

\begin{itemize}
  \item The manuscript defines a complete communication protocol rather than
  reporting sparse uplink traffic while ignoring the downlink needed to keep
  clients synchronized.
  \item The novelty boundary is careful. In particular, the paper distinguishes
  its communication-layer use of leaky integrate-and-fire ideas from federated
  training of spiking neural networks.
  \item The separation between trigger threshold and applied server quantum is
  conceptually useful and is explained clearly.
  \item Claims are appropriately scoped. The paper does not claim unconditional
  convergence, statistical significance, or measured hardware-energy savings.
  \item The development and held-out partitions are separated, paired streams
  are used across methods, and the repository provides strong evidence and
  reproducibility checks.
\end{itemize}

\section*{Major concerns}

\begin{enumerate}
  \item \textbf{The primary comparisons are underpowered.}
  The headline benchmark means and standard deviations use three held-out
  realizations. Several differences are on the scale of the observed
  variability. For example, the CIFAR-10 accuracy gap to the strongest Strom
  point is only 0.38 percentage points, while the reported standard deviations
  are 1.12 and 0.49 points. The manuscript correctly avoids a significance
  claim, but still uses language such as ``best mean accuracy'' and
  ``nondominated'' to carry much of the empirical case. With only three paired
  observations, the stability of the ranking cannot be assessed credibly.

  The paper should report per-seed paired differences, uncertainty intervals
  for those differences, and a predeclared paired analysis. Extending the main
  comparisons to at least ten held-out seeds would be the cleanest remedy. The
  ten-seed alignment factorial does not replace this because it studies a
  diagnostic on one Fashion-MNIST operating point rather than the headline
  method comparisons.

  \item \textbf{The experiments do not isolate the mechanism that supports the
  novelty claim.}
  The leakage ablation reports essentially no difference between
  $\rho=0.999$ and $\rho=1$. The coupled-resolution ablation sets
  $q_r=\vartheta=0.025$ and performs very poorly, but that comparison changes
  both the magnitude and schedule of the applied update. It therefore does not
  by itself show that decoupling is the causal reason for the reported gain.
  The manuscript also lacks a matched comparison between full reset and
  conservative threshold subtraction with the same trigger, quantum schedule,
  tuning budget, and traffic target.

  A convincing mechanism study should cross at least leak versus no leak,
  full reset versus threshold subtraction, and coupled versus independently
  tuned quantum. The comparison should control traffic or present a Pareto
  curve. In the current results, the specifically neuromorphic feature that is
  easiest to identify, leakage, appears empirically inactive.

  \item \textbf{The baseline set and tuning parity need strengthening.}
  Strom is the correct closest predecessor, and error-feedback controls are
  useful. However, STC is discussed as an important bidirectional sparse
  compressor but is not evaluated. At least one stronger modern
  bidirectional or event-triggered FL control would make the novelty and
  performance case more credible. The paper should also give the complete
  search space and number of development trials for every method in the main
  text or a compact appendix artifact. It is currently difficult to determine
  whether Event-FedAvg and all controls received comparable tuning effort.

  \item \textbf{The empirical scope is small for the breadth of the story.}
  All experiments use ten clients with full participation. The CIFAR-10 model
  has only 20,570 parameters and reaches about 48\% accuracy. This is useful as
  a controlled experiment, but it does not yet demonstrate that the address
  overhead and event dynamics remain favorable for a standard model, a larger
  client population, or client subsampling. One targeted stress test with
  partial participation and a more standard CIFAR architecture would add more
  value than another small variation of the same setting.

  \item \textbf{The optimization result remains an interface condition rather
  than a convergence explanation.}
  The descent inequality is valid and the alignment decomposition is
  informative. However, the stationarity result assumes a positive conditional
  alignment inequality that already contains the central optimization
  difficulty. The empirical statistic $\widehat{\kappa}_{\mathcal A}$ measures
  realized finite-trajectory alignment and does not test the conditional
  expectation or bound $\beta_r$. As written, the result does not explain how
  the leak, threshold, reset rule, heterogeneity, and local-step count guarantee
  the required condition.

  The authors should either derive a non-vacuous sufficient condition for a
  restricted setting or present the result more explicitly as a diagnostic
  accounting framework rather than convergence theory. A full proof of the
  event-budget result and precise measurability assumptions should be included
  in an archival supplement if the venue permits one.
\end{enumerate}

\section*{Questions for the authors}

\begin{enumerate}
  \item How many hyperparameter configurations were evaluated for
  Event-FedAvg and for each baseline before the operating points were frozen?
  \item Does full reset still outperform Strom-style residual subtraction when
  both use the same $q_r$, threshold grid, and total development budget?
  \item How often does the downlink use sparse replay rather than a dense
  checkpoint, and how sensitive is total traffic to client count?
  \item Are the reported per-seed advantages paired consistently across every
  baseline, and what are the confidence intervals of those paired
  differences?
  \item What behavior is expected when only a subset of clients participates
  and client encoder states become stale for different lengths of time?
\end{enumerate}

\section*{Minor comments}

\begin{itemize}
  \item ``First passage'' may suggest a continuously monitored process. The
  implementation tests each coordinate once per round, so ``round-wise sampled
  threshold crossing'' would be more precise at first use.
  \item Algorithm~1 would be easier to audit if client and server state were
  separated explicitly and the packet/checkpoint decision were written as
  pseudocode rather than prose inside an enumerated list.
  \item State whether round identifiers, client identifiers, and any framing
  needed for ordered replay are included in the 64-bit packet header.
  \item Report checkpoint frequency and the uplink/downlink split for each
  headline operating point. A single total hides which part of the protocol
  produces the saving.
  \item Explain why division by $\eta_r$ in the encoder is desirable when local
  learning-rate schedules vary, since it changes which training quantities the
  threshold treats as comparable over time.
  \item The current qualification about worst-class CIFAR-10 accuracy is useful
  and should remain visible in the final version.
  \item Figure~1 carries many distinct ideas. A simpler visual hierarchy and a
  separate synchronization inset would improve readability.
\end{itemize}

\section*{Revision threshold for acceptance}

I would reconsider the paper positively if a revision accomplished the
following four items:

\begin{enumerate}
  \item extend the headline comparisons to a defensible number of paired seeds
  and report paired uncertainty.
  \item add a controlled reset/leak/quantum ablation that isolates the claimed
  operator semantics.
  \item add at least one strong bidirectional compression baseline, preferably
  STC, with transparent tuning parity.
  \item either strengthen the link between the encoder and the alignment
  condition or narrow the theoretical claim to a finite-trajectory diagnostic
  framework.
\end{enumerate}

Partial participation or a standard-scale CIFAR model would materially improve
the paper, but I view either one as a high-value strengthening experiment
rather than an absolute prerequisite if the four items above are addressed
convincingly within the venue's space constraints.

\medskip
\textbf{Confidential note to the area chair.}
This is a promising and unusually careful artifact. My negative recommendation
reflects insufficient evidence, not an identified correctness failure. I would
support discussion at the acceptance boundary if the authors provide stronger
paired results or if other reviewers verify the mechanism and novelty more
favorably.

\end{document}
